\section{Separate component order from byte order}\label{rgb-encoding}
For pixel index $i$, packed RGB888 places red at byte $3i$, green at $3i+1$ and blue
at $3i+2$. Each byte is interpreted unsigned, so a load must produce zero through
255. A signed byte load would turn values above 127 into negative integers unless
the code explicitly removed that sign extension.

The supplied byte declarations use signed initializer notation where the assembler
requires it. For example, \texttt{-9} stores the eight-bit pattern \texttt{f7}.
An unsigned byte load reads that pattern as 247. The initializer spelling does not
change the channel's unsigned interpretation.

RGB565 first defines a numerical halfword. Its fields are independent of how the
halfword is stored. This exercise then chooses little-endian storage: the low eight
bits appear at byte $2i$ and the high eight bits at $2i+1$. Unaligned slices make
individual byte loads and stores a suitable scalar baseline. A halfword load or
store has additional native alignment requirements; mapped memory alone does not
establish that its address is aligned.

\begin{figure}[htbp]
  \centering
  \begin{tikzpicture}[alt={The sixteen-bit RGB565 value is divided into three adjoining fields. Red occupies bits fifteen through eleven, green ten through five, and blue four through zero. Red and blue each use five bits; green uses six. A separate storage row shows low byte before high byte.},x=5.5mm,y=9mm]
    \draw (0,0) rectangle (5,1);
    \draw (5,0) rectangle (11,1);
    \draw (11,0) rectangle (16,1);
    \node at (2.5,0.5) {Red, 5};
    \node at (8,0.5) {Green, 6};
    \node at (13.5,0.5) {Blue, 5};
    \node[above,font=\small] at (0,1) {15};
    \node[above,font=\small] at (5,1) {10};
    \node[above,font=\small] at (11,1) {4};
    \node[above,font=\small] at (16,1) {0};
    \node[below,font=\small] at (4,-0.4) {Memory byte zero: low byte};
    \node[below,font=\small] at (12,-0.4) {Byte one: high byte};
  \end{tikzpicture}
  \caption{Numerical fields and chosen storage order. Bit positions do not name separate memory addresses.}
\end{figure}

The diagram places the most significant field on the left, as is usual when writing
a binary number. Memory still stores the low byte first in this exercise. Do not
read the drawing from left to right as an address map.
